% !TeX program = xelatex
% docs/PAPER.tex 의 전문(preamble) + 표제 + 목차.
% tools/md2tex.py 가 이 파일을 읽어 본문 앞에 붙인다. 직접 조판하지 않는다.
\documentclass[10pt,a4paper]{article}

\usepackage{amsmath,amssymb}
\usepackage{booktabs,longtable,array}
\usepackage[margin=2.2cm]{geometry}
\usepackage{graphicx}
\usepackage{xcolor}
\usepackage{fancyvrb}
\usepackage{enumitem}
\usepackage{caption}
\usepackage[hidelinks,unicode]{hyperref}

\usepackage{xeCJK}
\setCJKmainfont[BoldFont={NanumMyeongjoBold}]{NanumMyeongjo}
\setCJKsansfont[BoldFont={Malgun Gothic Bold}]{Malgun Gothic}
\setCJKmonofont[BoldFont={Malgun Gothic Bold}]{Malgun Gothic}
% \XeTeXlinebreaklocale "ko" 는 쓰지 않는다. 켜면 ICU 가 한글을 글자 단위로 끊을 수 있게 되고,
% 낱말 사이 공백이 0 까지 수축해 단어가 모두 붙어 보인다. xeCJK 기본 줄바꿈만 쓴다.
\xeCJKsetup{CJKspace=true}

\graphicspath{{../figs/}{figs/}{./}}
\captionsetup{font=small,labelfont=bf,labelsep=period}
\captionsetup[figure]{name={그림}}
\captionsetup[table]{name={표}}

\setlength{\parindent}{0pt}
\setlength{\parskip}{0.55em}
\setlength{\emergencystretch}{3em}
\renewcommand{\arraystretch}{1.15}
\sloppy

% 한글 서술을 담은 강조상자. 수식 \boxed 는 줄바꿈이 안 되므로 이것으로 대체한다.
\newcommand{\keybox}[1]{%
  \par\smallskip\noindent
  \begin{center}\fbox{\parbox{0.93\linewidth}{\centering #1}}\end{center}%
  \par\smallskip}

\hypersetup{
  pdftitle={쌍동형 예인선 USV 의 3자유도 조종 동역학 식별},
  pdfauthor={윤원근 (ANSL, 충남대학교 자율운항시스템공학과)}}

\begin{document}

\begin{center}
{\LARGE\bfseries 쌍동형 예인선 USV 의\\[2pt]
3자유도 조종 동역학 식별\par}
\vspace{6pt}
{\large 이론 \textperiodcentered{} 방법 \textperiodcentered{} 결과 통합본\par}
\vspace{14pt}
{\normalsize\bfseries Tug 1: 완전구동 쌍동형 무인수상정의\\[2pt]
출력오차 기반 3자유도 조종 모델 식별과 궤적 수준 검증\par}
\vspace{16pt}
{\normalsize 윤원근\par}
{\small AI 필드 로보틱스 연구실 (ANSL)\\
충남대학교 자율운항시스템공학과\par}
\vspace{10pt}
{\small 작성 2026-11-21 \textperiodcentered{} 저장소
\href{https://github.com/wkyouncnu/tug1-3dof-sysid}{\texttt{wkyouncnu/tug1-3dof-sysid}}\par}
\end{center}

\vspace{4pt}

\begin{quote}\small
\textbf{이 문서의 위치.} 본 문서는 \texttt{THEORY.md}(이론), \texttt{METHOD.md}(절차),
\texttt{RESULTS.md}(수치)를 \textbf{논문 체계로 통합한 단일 문서}다. 중복을 제거하고
수식을 LaTeX 로, 절차를 알고리즘으로 다시 썼다. 세 개별 문서는 세부 참조용으로 그대로 남겨 둔다.

\textbf{조판 경로 메모.} 본 문서의 PDF 는 \texttt{PAPER.tex} 를 \textbf{XeLaTeX + xeCJK}
(TinyTeX / TeX Live 2026) 로 직접 조판했다. 저장소의 \texttt{tools/md2html.py} + Chrome 인쇄
경로는 Markdown 을 HTML 로만 바꾸므로 \texttt{\$...\$} 수식을 조판하지 못한다. 한글 조판은
\texttt{NanumMyeongjo}(본문) / \texttt{Malgun Gothic}(산세리프 \textperiodcentered{} 고정폭)
으로 해결되어 KaTeX/MathJax 대체 경로는 사용하지 않았다. \texttt{PAPER.tex} 는
\texttt{tools/md2tex.py} 가 \texttt{docs/PAPER.md} 에서 생성한다.
\end{quote}

\vspace{6pt}
{\small\tableofcontents}
\clearpage
